\subsection{全skew torusから37パラメータ・六型への必要条件による有限化}
\label{sec:skew-finite-reduction}
本節は一般$q$での空性より強く，六次元skew torusの\emph{全域}で
H3包絡となり得るパラメータを有限集合に絞る．
表の各候補が実際にH3を含むという十分性までは主張しない．

\begin{lemma}[座標点のjetによる必要条件]
\label{lem:skew-coordinate-jets}
量子多項式環$B(q)$を包絡とするH3の三次関係$G$の単項式supportには，
各$i=0,1,2,3$について少なくとも一つの$2e_i+e_j$が含まれる．
$j=i$は$x_i^3$を意味する．
\end{lemma}
\begin{proof}
前節より$G\in W_1$であり，support内の指数の差は交換bicharacterのradicalに入る．
その差が生成する格子を$L\subset\mathbb Z^4$とし，
$\Gamma=\mathbb Z^4/L$と置く．交換bicharacterは$\Gamma$へ降りる．
有限生成Abel群の自由・巡回成分に基底を選び，三角積によって2-cocycleを作れば，
$B(q)$を$\Gamma$次数付き可換多項式環のcocycle twistとして表示できる．
巡回成分の次数による積の定義は，bicharacterの値の位数がその巡回次数を割るのでwell-definedである．
$G$はこの次数で斉次なので，$R$は同じsupportを持つ可換三次環
$A=k[x_0,x_1,x_2,x_3]/(f)$のtwistとなる．
係数は非零スカラー倍に変わり得るがsupportは変わらない．

$H=\operatorname{Spec}k[\Gamma]$は標数0で線型簡約な対角化可能群である．
$\Gamma$次数付き$qgr$のcocycle同値は，$H$同変な$qgr$のexact同値となる．
$H$が有限であるとはまだ仮定しない．
線型簡約性により同変Extのdescent spectral sequenceは退化し，
\[
 \Ext^n_H(M,N)=\Ext^n(M,N)^H
\]
となる．従って$\gldim\qgr R=2$から，その同変圏の大域次元も高々2である．

もし$i$について指定supportが一つもなければ，$f$と全ての一次偏微分は
第$i$座標点$P_i$で消える．$P_i$は$X=\operatorname{Proj}A$の特異点であり，
対角群$H$の固定点でもある．剰余体対象$k(P_i)$は同変化できる．
特異な可換局所環の剰余体の自己Extは任意に高い次数で非零であり，
支持が一点なのでこれは$X$上のExtと一致する．各Extは有限次元の$H$表現である．
非零weightに応じて標的の同変構造を逆指標で捩れば，
そのweightは同変Extの非零不変部分となる．
従って$(\operatorname{coh}X)^H$の大域次元は無限となり，矛盾する．
この議論は多重次数twistを無条件の$qgr$同値と取り違えていない．
\end{proof}

\begin{theorem}[skew H3包絡の有限候補リスト]
\label{thm:skew-finite-37}
六次元skew torus内でH3の二次包絡となる$q$は，
全係数が2乗根であるか，全係数が3乗根であるかのいずれかである．
座標表示を固定すると，可能性は高々37パラメータであり，
座標置換の範囲では次の六候補型に限られる．
$q_{ij}=\zeta_N^{e_{ij}}$，指数の順序を$01,02,03,12,13,23$とする．
\[
\begin{array}{c|c|r|c|r}
 N&(e_{01},e_{02},e_{03},e_{12},e_{13},e_{23})
 &\text{軌道長}&\text{CY}&\dim W_1\\\hline
 1&(0,0,0,0,0,0)&1&\text{yes}&20\\
 2&(0,0,0,1,1,1)&4&\text{yes}&5\\
 2&(0,1,1,1,1,1)&6&\text{no}&8\\
 3&(0,0,0,1,2,1)&8&\text{yes}&5\\
 3&(0,1,2,1,2,2)&12&\text{yes}&6\\
 3&(0,1,2,2,1,0)&6&\text{yes}&4
\end{array}
\]
表は必要条件の完全な有限絞り込みであり，
全行のrank-one条件・全係数の滑らかさ・成分性の十分性を主張しない．
非CYの行の全admissible係数と幾何学的同型群oidは
定理\ref{thm:fixed-sign-classification}で分類した．
\end{theorem}
\begin{proof}
補題で各$i$に対して一つの$j=f(i)$を選ぶ．$f$は4元集合の自己写像であり，
全部で$4^4=256$通りである．$a_i=2e_i+e_{f(i)}$，$d_i=a_i-\mathbf1$と置く．
$G\in W_1$から，全$i$と全$h$に対して
\[
 \prod_{j=0}^3 q_{jh}^{(d_i)_j}=1.
\]
六つの$q_{ab}$ ($a<b$)の指数関係行列$M_f$は16行6列で，
その$(i,h)$行，$(a,b)$列の成分は
\[
 (d_i)_a\,\delta_{bh}-(d_i)_b\,\delta_{ah}.
\]
座標置換で$f$を共役して，19個の自己写像型を調べればよい．
整数Smith標準形は次の五種類で尽くされる．
\[
\begin{array}{c|r}
\text{Smith対角の絶対値}&256通りのうちの個数\\\hline
(1,1,1,1,1,1)&149\\
(1,1,1,1,1,2)&84\\
(1,1,1,1,1,3)&18\\
(1,1,1,2,2,2)&4\\
(1,1,1,3,3,3)&1
\end{array}
\]
この整数計算の証明書を
\texttt{computations/cubic\_skew\_jet\_lattices.json}に保持する．
各19型の$16\times16$および$6\times6$整数行列$U,V$について
$UM_fV=D$，$\det U=\det V=\pm1$を独立に有理数演算で検査する．
従ってパラメータトーラスの解群はSmith対角で与えられ，
各$q_{ab}$の位数は同じ$N\in\{1,2,3\}$を割る．
ここは有限標数の消滅から特性零を推論した箇所ではなく，整数格子の恒等式である．

続いて$2^6=64$個と$3^6=729$個のroot指数を全列挙する．
$W_1$が座標jet条件を満たすのは，前者で23個，後者で27個である．
恒等パラメータが重複するため総数は49個となる．
符号型23個のうち12個では$W_1$の全単項式が同じ変数で割り切れる．
この場合の可換untwistの三次多項式は常に可約であり，
一次因子と二次因子の射影交点で特異となる．
$q$が有限位数になったので，support差で割った有限群$\Gamma$でcocycleを取れる．
前定理の有限同変比較により，その商圏は大域次元2になり得ない．
従って12個を除外でき，$49-12=37$個が残る．
残りのroot指数の座標置換軌道と$W_1$を全列挙すると表示の六行を得る．
整数変換・全列挙・共通因子の排除を
\texttt{python3 computations/verify\_cubic\_skew\_jets.py}が検証する．
デフォルトの検証器はPython標準ライブラリのみを使い，
整数標準形の生成に用いたSymPyへ依存せず証明書の全等式を再検査する．
\end{proof}

この有限化は，単なる27個のFermatパラメータの列挙と異なり，
任意の非中心三次元・非CY ambientを許した六次元skew torusを出発点とする．
各候補型のadmissible係数と内在的rankは次節で判定する．
skewでないAS4 ambientの被覆は，この有限化の範囲外である．
